% Source-derived chapter 07: Geometrically rational elliptic surfaces
% Original source: no-enriques-surface-over-the-integers
\chapter{Geometrically rational elliptic surfaces}
\label{no-enriques-surface-over-the-integers-chapter-07}
\source{no-enriques-surface-over-the-integers}

\begin{remark}[Source section]
\label{no-enriques-surface-over-the-integers-chapter-07-source}
\source{no-enriques-surface-over-the-integers}
\source{no-enriques-surface-over-the-integers}
This chapter reproduces the corresponding section of the retrieved
LaTeX source, including its definitions, statements, examples, and
proofs. It has not yet been translated into Lean.
\notready
\end{remark}

% The source section title was: Geometrically rational elliptic surfaces
\mylabel{Geometrically rational}

In this section we discuss the relation between Enriques surfaces and geometrically rational surfaces.
The results are well-known, but I could not find suitable references in the required generality.
Let $k$ be an arbitrary ground field of arbitrary characteristic $p\geq 0$.
A proper surface $J$ is called a \emph{geometrically rational surface} if it is   geometrically integral and the
base-change to $k^\alg$ is birational to the projective plane. Clearly we have $h^0(\O_J)=1$.

Suppose $J$ is such a surface that is  endowed with a   genus-one fibration $\phi:J\ra \PP^1$.
We assume that the fibration is jacobian and relatively minimal, and that the total space is smooth.
Fix a section $E\subset J$, and let $J\ra Z$ be the contraction of all vertical curves disjoint from the zero-section $E\subset J$.
The normal proper surface $Z$ is called the \emph{Weierstra\ss{} model}.
Write $\PP^1=\Spec k[t]\cup \Spec k[t^{-1}]$ for some indeterminate $t$.


\begin{proposition}
\source{no-enriques-surface-over-the-integers}
\notready
\mylabel{global weierstrass equations}
There are polynomials $a_i\in k[t]$ of degree $\deg(a_i)\leq i$
such that 
\begin{equation}
\label{two equations}
\source{no-enriques-surface-over-the-integers}
\begin{gathered}
y^2+a_1xy+a_3y=x^3+a_2x^2+a_4x+a_6,\\
y^2+a'_1xy+a'_3y=x^3+a'_2x^2+a'_4x+a'_6\qquad \text{with $a_i'=a_i/t^i$}
\end{gathered}
\end{equation}
are  Weierstra\ss{} equations for  $Z$ over the affine open sets $\PP^1\smallsetminus\{\infty\}=\Spec k[t]$
and $\PP^1\smallsetminus \{0\}=\Spec k[t^\inv]$, respectively.  
\end{proposition}

\proof
To simplify notation we write $\O_X(n)$ for the pullbacks of the invertible sheaves $\O_{\PP^1}(n)$.
The Canonical Bundle Formula \cite{Bombieri; Mumford 1977} gives $\omega_J=\O_X(d)$ for some integer $d$.
Base-changing to the algebraic closure $k^\alg$ we can apply \cite{Cossec; Dolgachev 1989}, Proposition 5.6.1 and conclude $d=-1$.
The Adjunction formula for $E\subset J$ gives the self-intersection number $(E\cdot E)=-1$,
and we have $\omega_{J/\PP^1}=\O_X(1)$.

As explained in \cite{Deligne 1975}, Section 1 the sheaves $\shE_i=f_*\O_X(iE)$ are locally free,
with $\rank(\shE_i)=i$ for  $i\geq 1$, and $\shE_0=\O_{\PP^1}$. The canonical inclusions of $\shE_i$
define an increasing filtration on $\shE=\shE_3$ with $\gr^*(\shE)=\bigoplus_{i=0,2,3}\O_{\PP^1}(-i)$.
The invertible sheaf $\shL=\O_X(3E)$ is relatively very ample, and yields a closed embedding
$Z\subset\PP(\shE)$.
On the projective line, we actually may choose splittings $\shE=\O_{\PP^1} \oplus\O_{\PP^1}(-2)\oplus\O_{\PP^1}(-3)$.

Write $\PP^1=\Proj k[t_0,t_1]$, such that $t=t_0/t_1$, and consider the global section $\omega=t_1$
of the invertible sheaf
$\O_{\PP^1}(1)=\Omega^1_{J/\PP^1}|E$. Over the affine open set $U=\Spec k[t]$, we choose local sections $x=t_0^{-2}$
and $y=t_0^{-3}$ in the second and third summand of $\shE$. 
As explained in loc.\ cit., there are  polynomials $a_i\in k[t]$ such that 
the first equation in \eqref{two equations} holds in the sheaf of graded rings $\Sym^\bullet(\shE)$ over $U$.
The situation over $U'=\Spec k[t^{-1}]$ is similar: Here we choose $\omega'=t_0$ and $x'=t_1^{-2} $ and $y'=t_1^{-3}$.
This gives polynomials $a_i'\in k[t^{-1}]$ with 
$y'^2+a_1'x'y'+\ldots = x'^3+\ldots$. Now we use that the  coefficients are unique, once the local sections 
in the summands of $\shE$ are chosen.
Multiplying the equation over $U$ with $t^{-6}=(t_0/t_1)^6$ and comparing coefficients with the equation over $U'$
we get $a_i'=a_i/t^i$. This gives the second equation in \eqref{two equations}, and  ensures $\deg(a_i)\leq i$.
\qed

\medskip
Note that at least one of the  coefficients $a_i$ in  \eqref{two equations} is non-constant, and at least one is not divisible by $t^i$,
because otherwise the geometrically rational smooth surface $J$ would be isomorphic to the product of $\PP^1$ and some elliptic curve.
Furthermore, this property of the coefficients remains  true after any degree-preserving change of coordinates.

Now let $Y$ be an Enriques surface, and suppose there is a genus-one fibration $\varphi:Y\ra\PP^1$.
Let $\eta\in\PP^1$ be the generic point, with function field $K=k(t)$. 
If the fibration is elliptic, $J_\eta=\Pic^0_{Y_\eta/K}$ is an elliptic curve.
If the fibration is quasielliptic, $\Pic^0_{Y_\eta/K}$ is a twisted form of the additive group $\GG_{a,K}$,
and we write $J_\eta$ for its unique regular compactification.
In both cases, $J_\eta$ is a regular curve with $h^0(\O_{J_\eta})=h^1(\O_{J_\eta})=1$, and we write
$\phi:J\ra\PP^1$ for the resulting relatively minimal genus-one fibration.
It comes with a zero-section,  and is thus a jacobian fibration.

\begin{proposition}
\source{no-enriques-surface-over-the-integers}
\notready
\mylabel{jacobian geometrically rational}
The regular surface $J$ is geometrically rational.
\end{proposition}

\proof
Let $U\subset Y$ be the open set where the coherent sheaf $\Omega^1_{Y/\PP^1}$ is invertible.
It is dense, because the generic fiber $Y_\eta$ is geometrically reduced, hence the image
$V=\varphi(U)$ is a dense open set. 
For each point $b\in V$, the fibration $\varphi:Y\ra\PP^1$ acquires a section after base-changing
to the strict henselization  $\O_{\PP^1,b}\subset R$, and we get $Y\times_{\PP^1}\Spec(R)\simeq J\times_{\PP^1}\Spec(R)$.
It follows that the scheme $J$ is smooth on the open set $\psi^{-1}(V)$, and we infer that the base-change $\bJ$ remains integral.

Moreover, the jacobian for the base-change $\bY$ coincides with the base-change of the jacobian $\bJ$, at least over some
dense open set in $\PP^1_{k^\alg}$, so we may assume from the start that $k$ is algebraically closed.
The assertion now follows from \cite{Cossec; Dolgachev 1989}, Theorem 5.7.2.
\qed

\medskip
In turn, the genus-one fibration $\varphi:Y\ra \PP^1$ yields polynomials $a_i\in k[t]$ of degree $\deg(a_i)\leq i$
such that the two Weierstra\ss{} equations in \eqref{two equations} describe the Weierstra\ss{} model $Z\ra \PP^1$
of the jacobian fibration $\phi:Y\ra\PP^1$. We shall see that for $k=\FF_2$, the possibilities for these polynomials
impose
strong restrictions on the Kodaira symbols occurring for $\varphi$.


%===========================================================
